\documentclass[11pt,a4paper]{scrartcl}
\usepackage[ngerman]{babel}
\usepackage[a4paper,margin=24mm]{geometry}
\usepackage{microtype,amsmath,amssymb,booktabs,longtable,hyperref}
\hypersetup{colorlinks=true,pdftitle={Prä-raumzeitliche Ontologie},pdfauthor={Ingolf Lohmann}}
\title{Prä-raumzeitliche Ontologie der Unterscheidung\\
\large Unterschied, Relation, Information, emergente Raumzeit und epistemische Grenzen}
\author{Ingolf Lohmann}
\date{4. August 2026}
\begin{document}
\maketitle
\begin{abstract}
Dieses Dokument formuliert eine prä-raumzeitliche Ontologie, ohne einen leeren Raum,
eine Uhr oder eine zeitliche Entstehungsfolge vorauszusetzen. ``Prä-raumzeitlich''
bezeichnet eine Ordnung der Erklärungsabhängigkeit, kein zeitliches Davor. Ausgangspunkt
ist Unterscheidbarkeit als minimaler relationaler Sachverhalt. Operational wirksame
Unterschiede werden als Information behandelt; aus höherstufigen Relationen, Quantität,
Ordnung, Invarianz und gerichteter Wirksamkeit kann eine Raumzeitbeschreibung emergieren.
Formale Aussagen, empirische Befunde, Modelle, offene Fragen, normative Setzungen und
Glaubensaussagen werden strikt getrennt.
\end{abstract}

\section{Epistemischer Vertrag}
Die zentrale Typregel lautet:
\[
\text{formaler Satz}\neq\text{empirischer Befund}\neq\text{Modell}
\neq\text{offene Frage}\neq\text{Norm}\neq\text{Glaube}.
\]
Information ist nicht automatisch Wahrheit; Korrelation und Sequenz sind nicht
automatisch Kausalität; ein Erlebnis ist nicht bereits der Beweis seiner bevorzugten
Deutung; ``nicht widerlegt'' bedeutet nicht ``bewiesen''.

\section{Prä-raumzeitliche Grundbegriffe}
\subsection{Unterscheidbarkeit}
Jede Aussage, Messung oder Identifikation setzt voraus, dass mindestens zwei mögliche
Zustände unterscheidbar sind. Das absolute Nichts kann nicht als gewöhnlicher Zustand
behandelt werden, weil bereits seine Bezeichnung eine Unterscheidung einführt.
Dies ist eine methodische Minimalbedingung, kein vollständiger kosmologischer Beweis.

\subsection{Information}
Ein Unterschied wird operational zu Information, wenn er registriert, stabilisiert,
übertragen, rekonstruiert oder für weitere Zustände folgenreich werden kann. Diese
Definition ist mit der mathematischen Kommunikationstheorie vereinbar
\citep{Shannon1948a,Shannon1948b}, sagt aber allein nichts über Wahrheit oder Bedeutung.

\subsection{Relation, Quantität und Ordnung}
Der Unterschied ist bereits eine minimale Nichtidentitätsrelation. Höherstufige
Relationen verbinden unterscheidbare Zustände; Quantität erlaubt Anzahl, Verhältnis
und Maß; Ordnung legt zulässige Verknüpfungen, Invarianten und Ausschlüsse fest.

\subsection{Kausalität}
Kausalität ist eine besondere gerichtete Relation. In interventioneller Schreibweise
wird eine Abhängigkeit relevant, wenn
\[
P(Y\mid \operatorname{do}(X=x))
\]
sich mit $x$ ändert. Kausalität ist daher nicht bloß zeitliche Reihenfolge
\citep{Pearl1995}.

\section{Emergenz der Raumzeit}
Eine zeitliche Ordnung kann aus einer konsistenten Kausalordnung und vergleichbaren
Veränderungen hervorgehen. Eine räumliche Ordnung kann aus stabiler Trennung,
Nachbarschaft, Erreichbarkeit und Maß hervorgehen. Raumzeit ist dann nicht zunächst
ein leerer Behälter, sondern die geometrische Beschreibung stabilisierter
Kausal- und Trennungsrelationen.

Mathematische Resultate zeigen, dass unter geeigneten Bedingungen die kausale Struktur
wesentliche geometrische Information einer Raumzeit bestimmt \citep{Malament1977}.
Causal-Set-Ansätze modellieren Raumzeit aus diskreter Kausalordnung
\citep{Bombelli1987}; thermodynamische und verschränkungsbasierte Programme untersuchen
weitere Emergenzmechanismen \citep{Jacobson1995,VanRaamsdonk2010}. Daraus folgt jedoch
keine einzig bewiesene Fundamentalontologie.

Die Abhängigkeit wird deshalb nur typisiert:
\[
\text{Unterschied}\Rightarrow\text{Information}\Rightarrow
\text{Relation/Quantität}\Rightarrow\text{Ordnung/Invarianz}\Rightarrow
\text{gerichtete Wirksamkeit}\Rightarrow
\text{Kausalordnung + Trennung + Maß}\Rightarrow\text{Raumzeitmodell}.
\]
Jeder Pfeil benötigt eine eigene Definition, ein Modell und gegebenenfalls empirische
Prüfung.

\section{Physische, biologische und kognitive Emergenz}
Raumzeit allein erzeugt keine konkrete Welt. Zusätzliche Dynamik beschreibt Felder,
Wellen, Materie und Erhaltung. Leben erfordert selbsterhaltende Organisation,
Informationsspeicherung, Fehlerkorrektur und offene Stoff- und Energieflüsse.
Kognition verbindet Wahrnehmung, Gedächtnis, Modellbildung und Handlungswahl.
Bewusstsein bleibt eine offene Theoriefrage; beim Menschen besteht starke Evidenz
für eine enge Abhängigkeit von organisierter Hirndynamik
\citep{Boveroux2010,Krom2020}.

\section{Wissen, Glauben und Jenseitsfragen}
Wissenschaftliches Wissen verlangt klare Aussagen, öffentlich prüfbare Evidenz,
Fehlerkontrollen, unabhängige Prüfung und Korrekturfähigkeit. Glaube kann Vermutung,
Vertrauen, Hoffnung oder religiöse Bindung bedeuten. Er ist nicht automatisch
irrational, besitzt aber ohne prüfbare Evidenz nicht denselben Status wie
wissenschaftliches Wissen.

Die Ontologie beweist Gott weder noch widerlegt sie Gott. Sobald eine konkrete
Glaubensbehauptung reproduzierbare Beobachtungsfolgen beansprucht, wird genau diese
Folge empirisch prüfbar. Bleibt sie definitionsgemäß ohne Beobachtungsfolge, ist sie
keine naturwissenschaftliche Hypothese.

Nahtoderfahrungen sind reale Berichte von Menschen in reversiblen Grenz- und
Reanimationssituationen \citep{Parnia2023}. Sie beweisen nicht den Fortbestand
personaler Bewusstheit nach dauerhaft festgestelltem Ausfall der gesamten Hirnfunktion;
dessen Diagnose besitzt eigene medizinische Kriterien
\citep{Greer2020,Greer2023}. Ein persönliches Jenseits ist damit wissenschaftlich
nicht etabliert, aber auch nicht durch Spott logisch widerlegt.

\section{Telepathie und Teletubbies}
Teletubbies sind Figuren einer Kindersendung. Telepathie ist die Behauptung einer
Übertragung neuer, spezifischer Information ohne bekannten Kommunikationskanal.
Ein gemeinsamer Wortanfang ist keine Sachidentität. Quantenverschränkung erlaubt
keinen kontrollierbaren überlichtschnellen Nachrichtenkanal
\citep{Ghirardi1982}. Die parapsychologische Literatur enthält umstrittene positive
Meta-Analysen \citep{TressoldiStorm2024}; unabhängige Replikationen liefern keinen
robust etablierten allgemeinen Telepathiekanal \citep{StedallTressoldi2025}.
Der korrekte Status lautet: prüfbare Behauptung, kein gesichertes Wissen.

\section{Verantwortung und Ethik}
Aus deskriptiven Tatsachen folgt ohne normative Brückenregel kein konkretes Sollen.
Verantwortung wächst plausibel mit kausaler Reichweite, Erkenntnismöglichkeit und
Alternativenkontrolle. Wissenschaft kann als rekursiver Verantwortungsoperator auf
Relationen verstanden werden: Sie prüft Unterschiede, Abhängigkeiten und Wirkungen,
korrigiert Modelle und macht Folgen sichtbar.

\section{Falsifizierbares Forschungsprogramm}
Eine physikalische Theorie muss Mikrozustände, Dynamik, Grobkörnigkeit, Dimension,
Lorentzstruktur, Metrik, bekannte Grenzfälle und neue Vorhersagen liefern. Telepathie
würde präregistrierte Multi-Labor-Studien, sensorische Abschirmung, kryptografisch
gebundene Randomisierung, unabhängige Replikation und die Übertragung neuer spezifischer
Information verlangen. Ein empirischer Jenseitsnachweis müsste unabhängig verifizierbare
Information liefern, deren Erwerb durch bekannte sensorische, technische, mnestische
und statistische Wege ausgeschlossen ist.

\section{Normalform und Grenzen}
\[
\begin{aligned}
&\text{Unterschied / Urrelation}\\
&\Rightarrow \text{operationale Information}\\
&\Rightarrow \text{höherstufige Relation und Quantität}\\
&\Rightarrow \text{Ordnung und Invarianz}\\
&\Rightarrow \text{gerichtete Wirksamkeit / Kausalität}\\
&\Rightarrow \text{Kausalordnung + relationale Trennung + Maß}\\
&\Rightarrow \text{emergente Raumzeitbeschreibung}\\
&\Rightarrow \text{Dynamik, Materie, Leben, Kognition}\\
&\Rightarrow \text{Wissen, Macht, Verantwortung und Zukunft}.
\end{aligned}
\]
Diese Kette ist keine vollständige Deduktion der Welt aus einem Wort. Sie ist ein
Prüfprogramm für notwendige Begriffe und offene Brücken.

\section{Schluss}
Ohne Unterschied keine Information. Ohne Relation kein Zusammenhang. Ohne gerichtete
Wirksamkeit keine Kausalordnung. Ohne Kausalordnung, Maß und Invarianz keine
Raumzeitbeschreibung. Ohne Prüfung kein Wissen. Ohne Wirkung keine Verantwortung.

\begin{verbatim}
SCIENTIFIC_STATUS = MODEL_WITH_FORMAL_AND_EMPIRICAL_BOUNDARIES
PEER_REVIEW = NOT_PERFORMED
COMPLETE_QUANTUM_GRAVITY_DERIVATION = NOT_CLAIMED
PHYSICAL_RETROCAUSALITY = NOT_ESTABLISHED
PAST_EVENT_MUTATION = NOT_ESTABLISHED
PERSONAL_AFTERLIFE = NOT_ESTABLISHED
TELEPATHIC_INFORMATION_CHANNEL = NOT_ESTABLISHED
PASS = NOT_CLAIMED
FINAL_PASS = NOT_CLAIMED
EFFECT_ACK_DONE = NOT_CLAIMED
\end{verbatim}

\bibliographystyle{plainnat}
\bibliography{REFERENCES}
\end{document}
